\documentclass[12pt]{article}
\usepackage[T1]{fontenc}
\usepackage[utf8]{inputenc}
\usepackage{lmodern}
\usepackage{textcomp}
\usepackage{enumitem}
\usepackage{geometry}
\geometry{margin=1in}
\begin{document}
\begin{center}
Answer Key - Machine Learning - Undergraduate Level Assessment 5 \
\small (Answers are ordered exactly as in the question paper.)
\end{center}

\section*{Multiple Choice}
\begin{enumerate}[label=\arabic*.]
\item \textbf{Answer:} C
\\ \textit{Options:} A) 0; B) 1; C) 2; D) 3
\\ \textit{Note:} The dimensionality of the null space of the matrix A, which consists of three identical rows, is 2 because the rank of the matrix is 1, leading to a nullity of 3 - 1 = 2 according to the rank-nullity theorem.
\item \textbf{Answer:} A
\\ \textit{Options:} A) True, True; B) False, False; C) True, False; D) False, True
\\ \textit{Note:} Both statements are true because ResNets are indeed feedforward networks that utilize skip connections but do not employ self-attention mechanisms, while Transformers are also feedforward networks that specifically leverage self-attention for processing input data.
\item \textbf{Answer:} C
\\ \textit{Options:} A) Optimize a convex objective function; B) Can only be trained with stochastic gradient descent; C) Can use a mix of different activation functions; D) None of the above
\\ \textit{Note:} Neural networks can employ various activation functions across different layers to optimize learning and capture non-linear relationships in data, enhancing their performance and versatility.
\item \textbf{Answer:} B
\\ \textit{Options:} A) 0; B) 1; C) 2; D) 3
\\ \textit{Note:} The rank of the matrix A is 1 because all rows are linearly dependent and can be expressed as scalar multiples of each other, indicating that there is only one linearly independent row.
\item \textbf{Answer:} D
\\ \textit{Options:} A) True, True; B) False, False; C) True, False; D) False, True
\\ \textit{Note:} The correct answer is false for Statement 1 because while ReLU (Rectified Linear Unit) functions are indeed monotonic (increasing), sigmoid functions are also monotonic, and true for Statement 2 because neural networks trained with gradient descent do not guarantee convergence to the global optimum due to the presence of local minima and saddle points in non-convex loss landscapes.
\end{enumerate}

\section*{True / False}
\begin{enumerate}[label=\arabic*.]
\setcounter{enumi}{5}
\item \textbf{Answer:} True
\\ \textit{Options:} A) True; B) False
\\ \textit{Note:} As the number of training examples increases, the model gains more information about the underlying data distribution, which reduces the impact of noise and leads to lower variance in predictions.
\item \textbf{Answer:} False
\\ \textit{Options:} A) True; B) False
\\ \textit{Note:} Bayesians advocate for the use of prior distributions in statistical models to incorporate prior knowledge, while frequentists typically do not use prior distributions and rely on data alone for inference.
\item \textbf{Answer:} True
\\ \textit{Options:} A) True; B) False
\\ \textit{Note:} The statement is true because Naive Bayes assumes that the presence or absence of a particular attribute is independent of the presence or absence of any other attribute when the class value is known, which simplifies the computation of probabilities in the model.
\item \textbf{Answer:} False
\\ \textit{Options:} A) True; B) False
\\ \textit{Note:} The statement is false because a Bayesian network with four nodes (H, U, P, W) requires more than 7 independent parameters due to the need to account for the conditional dependencies and relationships among these nodes, which typically exceed that number.
\item \textbf{Answer:} False
\\ \textit{Options:} A) True; B) False
\\ \textit{Note:} High entropy in classification indicates a high level of disorder or uncertainty in the data, meaning that the partitions contain a mixture of class labels rather than being pure and consisting of a single class label.
\end{enumerate}

\section*{Long Form Response}
\begin{enumerate}[label=\arabic*.]
\setcounter{enumi}{10}
\item \textbf{Answer:} def is\_decimal(num): | return bool(result)
\\ \textit{Note:} correct because it outlines the creation of a function that utilizes regular expressions (regex) to validate whether a given decimal number conforms to the specified precision of two decimal places, effectively returning a boolean result indicating the validity of the input.
\item \textbf{Answer:} def count\_Rectangles(radius): | return rectangles
\\ \textit{Note:} correct as it suggests a function structure that can be further developed to calculate the number of rectangles that can fit within a circle of a given radius by utilizing geometric principles and properties of shapes.
\end{enumerate}

\vfill
\noindent\textit{End of Answer Key}
\end{document}